\documentclass{ximera}

\title{Exercises for Functions and their Representations} 
\license{CC BY-NC-SA 4.0}

\begin{document}

\begin{abstract}
\end{abstract}
\maketitle


%In Exercises \ref{mappingfirst} - \ref{mappinglast}, d




\begin{enumerate}

\item  The mapping $M$ is not a function since  `Tennant' is matched with both `Eleven' and `Twelve.'

\item  The mapping $C$ is a function since each input is matched with only one output.  The domain of $C$ is $\{ \text{Hartnell}, \text{Cushing}, \text{Hurndall}, \text{Troughton} \}$ and the range is $\{\text{One}, \text{Two} \}$.  We can represent $C$ as the following set of ordered pairs:  $\{ (\text{Hartnell}, \text{One}), (\text{Cushing}, \text{One}),  (\text{Hurndall}, \text{One}), (\text{Troughton}, \text{Two}) \}$

%\setcounter{HW}{\value{enumi}}
\end{enumerate}

\begin{enumerate}

%\setcounter{enumi}{\value{HW}}

\item  In this case, $y$ is a function of $x$ since each $x$ is matched with only one $y$.  

The domain is $\{ -3, -2, -1,0,1,2,3 \}$ and the range is $\{ 0,1,2,3 \}$.  

As ordered pairs, this function is $\{ (-3,3), (-2,2), (-1,1), (0,0), (1,1), (2,2), (3,3) \}$

\item  In this case, $y$ is not a function of $x$ since there are $x$ values matched with more than one $y$ value.  For instance, $1$ is matched both to $1$ and $-1$.

%\setcounter{HW}{\value{enumi}}
\end{enumerate}


\begin{enumerate}

%\setcounter{enumi}{\value{HW}}

\item    The mapping is a function since given any word, there is only one answer to `how many letters are in the word?'  For the range, we would need to know what the length of the longest word is and whether or not we could find words of all the lengths between $1$ (the length of the word `a') and it.  See \href{https://en.wikipedia.org/wiki/Longest_word_in_English}{\underline{here}}.

\item  Since Grover Cleveland was both the 22nd and 24th POTUS, neither mapping described in this exercise is a function.

\item  The outdoor temperature could never be the same for more than two different times - so, for example, it could always be getting warmer or it could always be getting colder.

%\setcounter{HW}{\value{enumi}}

\end{enumerate}


\begin{enumerate}
%\setcounter{enumi}{\value{HW}}

\item $f(2) = \frac{7}{4}$, $f(x) = \frac{2x+3}{4}$

\item $f(2) = \frac{5}{2}$, $f(x) = \frac{2(x+3)}{4} = \frac{x+3}{2}$  

%\setcounter{HW}{\value{enumi}}
\end{enumerate}


\begin{enumerate}
%\setcounter{enumi}{\value{HW}}

\item $f(2) = 7$, $f(x) = 2\left(\frac{x}{4} + 3\right) = \frac{1}{2} x + 6$   

\item $f(2) = \sqrt{7}$, $f(x) = \sqrt{2x+3}$ 

%\setcounter{HW}{\value{enumi}}
\end{enumerate}


\begin{enumerate}
%\setcounter{enumi}{\value{HW}}

\item $f(2) = \sqrt{10}$, $f(x) = \sqrt{2(x+3)} = \sqrt{2x+6}$

\item $f(2) = 2 \sqrt{5}$, $f(x) = 2\sqrt{x+3}$ 

%\setcounter{HW}{\value{enumi}}
\end{enumerate}


\begin{enumerate}
%\setcounter{enumi}{\value{HW}}

\item For $f(x) = 2x+1$ 


\begin{itemize}
\item $f(3) = 7$
\item $f(-1) = -1$
\item $f\left(\frac{3}{2} \right) = 4$
\end{itemize}


\begin{itemize}
\item  $f(4x) = 8x+1$
\item $4f(x) = 8x+4$
\item $f(-x) = -2x+1$
\end{itemize}


\begin{itemize}
\item  $f(x-4) = 2x-7$
\item $f(x) - 4 = 2x-3$
\item  $f\left(x^2\right) = 2x^2+1$
\end{itemize}


\item For $f(x) = 3-4x$ 


\begin{itemize}
\item $f(3) = -9$
\item $f(-1) = 7$
\item $f\left(\frac{3}{2} \right) = -3$
\end{itemize}


\begin{itemize}
\item  $f(4x) = 3-16x$
\item $4f(x) = 12-16x$
\item $f(-x) = 4x+3$
\end{itemize}


\begin{itemize}
\item  $f(x-4) = 19-4x$
\item $f(x) - 4 = -4x-1$
\item  $f\left(x^2\right) = 3-4x^2$
\end{itemize}




\item For $f(x) = 2 - x^2$ 


\begin{itemize}
\item $f(3) = -7$
\item $f(-1) = 1$
\item $f\left(\frac{3}{2} \right) = -\frac{1}{4}$
\end{itemize}


\begin{itemize}
\item  $f(4x) = 2-16x^2$
\item $4f(x) = 8-4x^2$
\item $f(-x) = 2-x^2$
\end{itemize}


\begin{itemize}
\item  $f(x-4) = -x^2+8x-14$
\item $f(x) - 4 = -x^{2} - 2$
\item  $f\left(x^2\right) = 2-x^4$
\end{itemize}


\item For $f(x) = x^2 - 3x + 2$ 


\begin{itemize}
\item $f(3) = 2$
\item $f(-1) = 6$
\item $f\left(\frac{3}{2} \right) = -\frac{1}{4}$
\end{itemize}


\begin{itemize}
\item  $f(4x) = 16x^2-12x+2$
\item $4f(x) = 4x^2-12x+8$
\item $f(-x) = x^2+3x+2$
\end{itemize}


\begin{itemize}
\item  $f(x-4) = x^2-11x+30$
\item $f(x) - 4 = x^2-3x-2$
\item  $f\left(x^2\right) = x^4-3x^2+2$
\end{itemize}




\item For $f(x) = 6$ 



\begin{itemize}
\item $f(3) = 6$
\item $f(-1) =6$
\item $f\left(\frac{3}{2} \right) = 6$
\end{itemize}


\begin{itemize}
\item  $f(4x) = 6$
\item $4f(x) = 24$
\item $f(-x) = 6$
\end{itemize}


\begin{itemize}

\item  $f(x-4) = 6$ 

\item $f(x) - 4 = 2$
     
\item  $f\left(x^2\right) = 6$

\end{itemize}




\item For $f(x) = 0$ 


\begin{itemize}
\item $f(3) = 0$
\item $f(-1) =0$
\item $f\left(\frac{3}{2} \right) = 0$
\end{itemize}


\begin{itemize}
\item  $f(4x) = 0$
\item $4f(x) = 0$
\item $f(-x) = 0$
\end{itemize}


\begin{itemize}

\item  $f(x-4) = 0$ 

\item $f(x) - 4 = -4$
     
\item  $f\left(x^2\right) = 0$

\end{itemize}


%\setcounter{HW}{\value{enumi}}
\end{enumerate}




\begin{enumerate}
%\setcounter{enumi}{\value{HW}}

\item For $f(x) = 2x-5$


\begin{itemize}

\item  $f(2) = -1$
\item  $f(-2) = -9$
\item  $f(2a) = 4a-5$

\end{itemize}

\begin{itemize}

\item  $2 f(a) = 4a-10$
\item $f(a+2) = 2a-1$
\item $f(a) + f(2) = 2a-6$

\end{itemize}


\begin{itemize}

\item  $f \left( \frac{2}{a} \right) = \frac{4}{a} - 5$ \\
$\hphantom{f \left( \frac{2}{a} \right)} = \frac{4-5a}{a}$

 

 

\item $\frac{f(a)}{2} =\frac{2a-5}{2}$

 

 


\item  $f(a + h) = 2a + 2h - 5$

\end{itemize}



\item For $f(x) = 5-2x$



\begin{itemize}

\item  $f(2) = 1$
\item  $f(-2) = 9$
\item  $f(2a) = 5-4a$

\end{itemize}


\begin{itemize}

\item  $2 f(a) = 10-4a$
\item $f(a+2) = 1-2a$
\item $f(a) + f(2) = 6-2a$

\end{itemize}


\begin{itemize}

\item  $f \left( \frac{2}{a} \right) = 5 - \frac{4}{a}$ \\
$\hphantom{f \left( \frac{2}{a} \right)} = \frac{5a-4}{a}$

 

 

\item $\frac{f(a)}{2} = \frac{5-2a}{2}$

 

 


\item  $f(a + h) = 5-2a-2h$

\end{itemize}




\item For $f(x) = 2x^2-1$


\begin{itemize}

\item  $f(2) = 7$
\item  $f(-2) = 7$
\item  $f(2a) = 8a^2-1$

\end{itemize}


\begin{itemize}

\item  $2 f(a) = 4a^2-2$
\item $f(a+2) = 2a^2+8a+7$
\item $f(a) + f(2) = 2a^2+6$

\end{itemize}


\begin{itemize}

\item  $f \left( \frac{2}{a} \right) = \frac{8}{a^2} - 1$ \\
$\hphantom{f \left( \frac{2}{a} \right)} = \frac{8-a^2}{a^2}$

 

 

\item $\frac{f(a)}{2} =  \frac{2a^2-1}{2}$

 

 


\item  $f(a + h) = 2a^2+4ah+2h^2-1$

\end{itemize}



\item For $f(x) = 3x^2+3x-2$


\begin{itemize}

\item  $f(2) = 16$
\item  $f(-2) = 4$
\item  $f(2a) = 12a^2+6a-2$

\end{itemize}


\begin{itemize}

\item  $2 f(a) = 6a^2+6a-4$
\item $f(a+2) = 3a^2+15a+16$
\item \small $f(a) + f(2) = 3a^2+3a+14$ \normalsize

\end{itemize}


\begin{itemize}

\item  $f \left( \frac{2}{a} \right) = \frac{12}{a^2} + \frac{6}{a} - 2$ \\
$\hphantom{f \left( \frac{2}{a} \right)} = \frac{12+6a-2a^2}{a^2}$

 

 

\item $\frac{f(a)}{2} =  \frac{3a^2+3a-2}{2}$

 

 


\item  $f(a + h) = 3a^2 + 6ah + 3h^2+3a+3h-2$

\end{itemize}



\item For $f(x) = 117$


\begin{itemize}

\item  $f(2) = 117$
\item  $f(-2) = 117$
\item  $f(2a) = 117$

\end{itemize}


\begin{itemize}

\item  $2 f(a) = 234$
\item $f(a+2) = 117$
\item $f(a) + f(2) = 234$

\end{itemize}


\begin{itemize}

\item  $f \left( \frac{2}{a} \right) = 117$ 

 

 

\item $\frac{f(a)}{2} = \frac{117}{2}$

 

 


\item  $f(a + h) = 117$

\end{itemize}




\item For $f(x) = \frac{x}{2}$


\begin{itemize}

\item  $f(2) = 1$
\item  $f(-2) = -1$
\item  $f(2a) = a$

\end{itemize}


\begin{itemize}

\item  $2 f(a) = a$

\item $f(a+2) = \frac{a+2}{2}$

 

 

\item $f(a) + f(2) = \frac{a}{2}+ 1$ \\
      $\hphantom{f(a) + f(2)} = \frac{a+2}{2}$

\end{itemize}


\begin{itemize}

\item  $f \left( \frac{2}{a} \right) = \frac{1}{a}$

 

 

\item $\frac{f(a)}{2} =  \frac{a}{4}$

 

 


\item  $f(a + h) = \frac{a+h}{2}$

\end{itemize}





%\setcounter{HW}{\value{enumi}}
\end{enumerate}

\begin{enumerate}
%\setcounter{enumi}{\value{HW}}

\item For $f(x) = 2x-1$,  $f(0) = -1$ and $f(x) = 0$ when $x = \frac{1}{2}$

\item For $f(x) =  3 - \frac{2}{5} x$, $f(0) = 3$ and $f(x) = 0$ when $x = \frac{15}{2}$

\item For $f(x) =  2x^2-6$, $f(0) = -6$ and $f(x) = 0$ when $x = \pm \sqrt{3}$

\item For $f(x) =  x^2-x-12$, $f(0) = -12$ and $f(x) = 0$ when $x = -3$ or $x=4$


%\setcounter{HW}{\value{enumi}}
\end{enumerate}


\begin{enumerate}
%\setcounter{enumi}{\value{HW}}


\item Function
\item Function
\item Function

%\setcounter{HW}{\value{enumi}}
\end{enumerate}


\begin{enumerate}
%\setcounter{enumi}{\value{HW}}


\item Not a function
\item Function
\item Not a function

%\setcounter{HW}{\value{enumi}}
\end{enumerate}


\begin{enumerate}
%\setcounter{enumi}{\value{HW}}

\item  Not a function
\item  Function
\item  Not a function

%\setcounter{HW}{\value{enumi}}
\end{enumerate}


\begin{enumerate}
%\setcounter{enumi}{\value{HW}}

\item   Function
\item   Not a function
\item Function

%\setcounter{HW}{\value{enumi}}
\end{enumerate}


\begin{enumerate}
%\setcounter{enumi}{\value{HW}}

\item Function
\item  Function
\item Not a function

%\setcounter{HW}{\value{enumi}}
\end{enumerate}


\begin{enumerate}
%\setcounter{enumi}{\value{HW}}

\item Function \\ domain = \{$-3$, $-2$, $-1$, $0$, $1$, $2$ ,$3$\} \\ range = \{$0$, $1$, $4$, $9$\}

 

 

\item Not a function

%\setcounter{HW}{\value{enumi}}
\end{enumerate}


\begin{enumerate}
%\setcounter{enumi}{\value{HW}}

\item  Function \\ domain = $\left\{ -7, -3, 3, 4, 5, 6 \right\}$ \\ range = $\left\{ 0,4,5,6,9 \right\}$


 

 

\item  Function \\ domain =   $\left\{ 1, 4, 9, 16, 25, 36, \ldots \right\} \\ = \left\{ x \, | \, \text{$x$ is a perfect square} \right\}$ \\ range =  $\left\{ 2, 4, 6, 8, 10, 12, \ldots \right\} \\ = \left\{ y \, | \, \text{$y$ is a positive even integer} \right\}$

%\setcounter{HW}{\value{enumi}}
\end{enumerate}


\begin{enumerate}
%\setcounter{enumi}{\value{HW}}

\item  Not a function

 

 

\item Function \\ domain $= \{x \, | \, \text{$x$ is irrational} \}$ \\ range $= \{ 1\}$

%\setcounter{HW}{\value{enumi}}
\end{enumerate}


\begin{enumerate}
%\setcounter{enumi}{\value{HW}}

\item Function \\ domain  $= \{x \, | \, 1, 2, 4, 8, \ldots \}$ \\ $= \{x \, | \, \text{$x = 2^{n}$ for some whole number $n$} \}$ \\ range $= \{ 0, 1, 2, 3, \ldots \}$ \\ $= \{y \, | \, \text{$y$ is any whole number}\}$

 

 

\item Function \\ domain $= \{x \, | \, \text{$x$ is any integer} \}$ \\ range $= \{y \, | \, \text{ $y$ is the square of an integer}\}$

%\setcounter{HW}{\value{enumi}}
\end{enumerate}


\begin{enumerate}
%\setcounter{enumi}{\value{HW}}

\item Not a function

 

 

\item Function \\ domain  $= \{x \, | \, -2 \leq x < 4 \} = [-2, 4)$, \\ range = \{$3$\}

%\setcounter{HW}{\value{enumi}}
\end{enumerate}


\begin{enumerate}
%\setcounter{enumi}{\value{HW}}


\item Function \\ domain $= \{x \, | \,  \text{$x$ is a real number} \} = (-\infty, \infty)$ \\  range $= \{y \, | \,  y \geq 0 \} = [0,\infty)$

 

 

\item  Not a function

%\setcounter{HW}{\value{enumi}}
\end{enumerate}


\begin{enumerate}
%\setcounter{enumi}{\value{HW}}

 \item \textbf{Horizontal Line Test:} A graph on the $xy$-plane represents $x$ as a function of $y$ if and only if no \textbf{horizontal} line intersects the graph more than once.
 
 %\setcounter{HW}{\value{enumi}}
\end{enumerate}





\begin{enumerate}
%\setcounter{enumi}{\value{HW}}

\item Function \\ domain = \{$-4$, $-3$, $-2$, $-1$, $0$, $1$\} \\ range = \{$-1$, $0$, $1$, $2$, $3$, $4$\}

 

 

\item Not a function

%\setcounter{HW}{\value{enumi}}
\end{enumerate}


\begin{enumerate}
%\setcounter{enumi}{\value{HW}}

\item Function \\ domain = $(-\infty, \infty)$ \\ range = $[1, \infty)$
 

 

\item Not a function 

%\setcounter{HW}{\value{enumi}}
\end{enumerate}


\begin{enumerate}
%\setcounter{enumi}{\value{HW}}

\item  

\begin{itemize}

\item Number \ref{graphfunctionfirstxy} represents $x$ as a function of $y$.

domain  = \{$-1$, $0$, $1$, $2$, $3$, $4$\} and range = \{$-4$, $-3$, $-2$, $-1$, $0$, $1$\}

\item Number \ref{graphfunctionlastxy} represents $x$ as a function of $y$.

domain = $(-\infty, \infty)$  and range = $[1, \infty)$


\end{itemize}

%\setcounter{HW}{\value{enumi}}
\end{enumerate}


\begin{enumerate}
%\setcounter{enumi}{\value{HW}}

\item Function \\ domain = $[2, \infty)$ \\ range = $[0, \infty)$

 

 

\item Function \\ domain = $(-\infty, \infty)$ \\ range = $(0, 4]$

%\setcounter{HW}{\value{enumi}}
\end{enumerate}



\begin{enumerate}
%\setcounter{enumi}{\value{HW}}

\item Not a function


 

 

\item Function \\ domain = $[-5,-3) \cup(-3, 3)$ \\ range = $(-2, -1) \cup [0, 4)$

%\setcounter{HW}{\value{enumi}}
\end{enumerate}


\begin{enumerate}
%\setcounter{enumi}{\value{HW}}

\item Only number \ref{graphfunctionfirstvw} represents $v$ as a function of $w$;  domain = $[0, \infty)$ and range = $[2, \infty)$

%\setcounter{HW}{\value{enumi}}
\end{enumerate}



\begin{enumerate}
%\setcounter{enumi}{\value{HW}}

\item  Function \\  domain =  $[-2, \infty)$ \\ range = $[-3, \infty)$

 

 

\item Not a function

%\setcounter{HW}{\value{enumi}}
\end{enumerate}


\begin{enumerate}
%\setcounter{enumi}{\value{HW}}

\item Function \\  domain =  $(-5, 4)$ \\ range = $(-4, 4)$

 

 

\item  Function \\ domain = $[0,3) \cup (3,6]$ \\ range = $(-4,-1] \cup [0,4]$

%\setcounter{HW}{\value{enumi}}
\end{enumerate}



\begin{enumerate}
%\setcounter{enumi}{\value{HW}}

\item None of numbers \ref{graphfunctionfirsttT} - \ref{graphfunctionlasttT}  represent $t$ as a function of $T$.

%\setcounter{HW}{\value{enumi}}
\end{enumerate}



\begin{enumerate}
%\setcounter{enumi}{\value{HW}}

\item  Function \\ domain = $(-\infty, \infty)$ \\ range = $(-\infty, 4]$

 

 

\item  Function \\ domain = $(-\infty, \infty)$ \\ range = $(-\infty, 4]$

%\setcounter{HW}{\value{enumi}}
\end{enumerate}



\begin{enumerate}
%\setcounter{enumi}{\value{HW}}

\item  Function \\  domain =  $[-2, \infty)$  \\ range =   $(-\infty, 3]$

 

 

\item  Function \\ domain = $(-\infty, \infty)$ \\ range = $(-\infty, \infty)$

%\setcounter{HW}{\value{enumi}}
\end{enumerate}



\begin{enumerate}
%\setcounter{enumi}{\value{HW}}

\item Only number \ref{graphfunctionfirstHs3} represents $s$ as a function of $H$;  domain =  $(-\infty, 3]$  and range =   $[-2, \infty)$


%\setcounter{HW}{\value{enumi}}
\end{enumerate}



\begin{enumerate}
%\setcounter{enumi}{\value{HW}}

\item  Function \\  domain =  $(-\infty, 0] \cup (1, \infty)$ \\ range =  $(-\infty, 1] \cup \{ 2\}$

 

 

\item  Function \\  domain =  $[-3,3]$ \\ range =  $[-2,2]$

%\setcounter{HW}{\value{enumi}}
\end{enumerate}



\begin{enumerate}
%\setcounter{enumi}{\value{HW}}

\item   Not a function

 

 

\item  Function \\ domain = $(-\infty, \infty)$ \\ range = $\{2\}$

%\setcounter{HW}{\value{enumi}}
\end{enumerate}



\begin{enumerate}
%\setcounter{enumi}{\value{HW}}

\item  Only number \ref{graphfunctionfirstut3} represents $t$ as a function of $u$;  domain = $(-\infty, \infty)$ and range=$\{2 \}$.

%\setcounter{HW}{\value{enumi}}
\end{enumerate}





\begin{enumerate}

%\setcounter{enumi}{\value{HW}}

\item  $f(-2) = 2$

\item $g(-2) = -5$

\item $f(2) = 3$

\item  $g(2) = 3$


%\setcounter{HW}{\value{enumi}}

\end{enumerate}




\begin{enumerate}

%\setcounter{enumi}{\value{HW}}

\item $f(0) = -1$

\item $g(0) = 0$

\item  $x  = -4, -1, 1$

\item  $t = -4, 0, 4$

%\setcounter{HW}{\value{enumi}}

\end{enumerate}




\begin{enumerate}

%\setcounter{enumi}{\value{HW}}

\item  Domain:  $[-5,3]$,  Range:  $[-5,4]$.

\item  Domain:  $[-4,4]$,  Range:  $[-5,5)$.

%\setcounter{HW}{\value{enumi}}

\end{enumerate}





\begin{enumerate}

%\setcounter{enumi}{\value{HW}}

\item 

$f(x) =2-x$

Domain: $(-\infty, \infty)$ 

Range:  $(-\infty, \infty)$


% \begin{mfpic}[15]{-3}{4}{-2}{4}

% \axes
% \tlabel[cc](4,-0.5){\scriptsize $x$}
% \tlabel[cc](0.5,4){\scriptsize $y$}
% \xmarks{-2,-1,1,2,3}
% \ymarks{-1,1,2,3}
% \tlpointsep{4pt}
% \tiny 
% \axislabels {x}{{$-2 \hspace{6pt}$} -2,{$-1 \hspace{6pt}$} -1, {$1$} 1, {$2$} 2, {$3$} 3}
% \axislabels {y}{{$-1$} -1, {$1$} 1, {$2$} 2, {$3$} 3}
% \normalsize
% \penwd{1.25pt}
% \arrow \reverse \arrow \function{-1.75,3.75, 0.1}{2-x}
% \point[4pt]{(-1,3), (0, 2), (1,1), (2,0), (3,-1)}
% \end{mfpic}


\item  

$g(t) = \dfrac{t - 2}{3}$

Domain: $(-\infty, \infty)$ 

Range: $(-\infty, \infty)$  

 

%  

% \begin{mfpic}[15]{-2}{5}{-2}{2}
% \axes
% \tlabel[cc](5,-0.5){\scriptsize $t$}
% \tlabel[cc](0.5,2){\scriptsize $y$}
% \xmarks{-1,1,2,3,4}
% \ymarks{-1,1}
% \tlpointsep{4pt}
% \tiny 
% \axislabels {x}{{$-1 \hspace{6pt}$} -1, {$1$} 1, {$2$} 2, {$3$} 3, {$4$} 4}
% \axislabels {y}{{$-1$} -1, {$1$} 1}
% \normalsize
% \penwd{1.25pt}
% \arrow \reverse \arrow \function{-2,5, 0.1}{(x - 2)/3}
% \point[4pt]{(-1,-1), (0, -0.6667), (1,-0.3333), (2,0), (3, 0.3333)}
% \end{mfpic}



\item 

$h(s) = s^2+1$

Domain: $(-\infty, \infty)$ 

Range:  $[1, \infty)$

 

 


% \begin{mfpic}[15]{-3}{3}{-1}{6}
% \axes
% \tlabel[cc](3,-0.5){\scriptsize $s$}
% \tlabel[cc](0.5,5.75){\scriptsize $y$}
% \xmarks{-2,-1,1,2}
% \ymarks{1,2,3,4,5}
% \tlpointsep{4pt}
% \axislabels {x}{{\tiny $-2 \hspace{8pt}$} -2, {\tiny $-1 \hspace{8pt}$} -1, {\tiny $1$} 1, {\tiny $2$} 2}
% \axislabels {y}{{\tiny $1$} 1, {\tiny $2$} 2, {\tiny $3$} 3, {\tiny $4$} 4, {\tiny $5$} 5}
% \penwd{1.25pt}
% \arrow \reverse \arrow \function{-2.1, 2.1, 0.1}{x**2+1}
% \point[4pt]{(-1,2), (0,1), (1,2), (2,5), (-2,5)}
% \end{mfpic}



\item 


$f(x) = 4-x^2$

Domain: $(-\infty, \infty)$ 

Range:  $(-\infty, 4]$

 

 


% \begin{mfpic}[15]{-3}{3}{-1}{5}
% \axes
% \tlabel[cc](3,-0.5){\scriptsize $x$}
% \tlabel[cc](0.5,5){\scriptsize $y$}
% \xmarks{-2,-1,1,2}
% \ymarks{1,2,3,4}
% \tlpointsep{4pt}
% \axislabels {x}{{\tiny $-2 \hspace{6pt}$} -2, {\tiny $-1 \hspace{6pt}$} -1, {\tiny $1$} 1, {\tiny $2$} 2}
% \axislabels {y}{{\tiny $1$} 1, {\tiny $2$} 2, {\tiny $3$} 3, {\tiny $4$} 4}
% \penwd{1.25pt}
% \arrow \reverse \arrow \function{-2.25,2.25,0.1}{4-(x**2)}
% \point[4pt]{(-1,3), (0,4), (1,3), (-2,0), (2,0)}
% \end{mfpic} 



\item   

$g(t) = 2$

Domain: $(-\infty, \infty)$ 

Range:  $\{2\}$

 

 


% \begin{mfpic}[15]{-3}{3}{-1}{4}
% \axes
% \tlabel[cc](3,-0.5){\scriptsize $t$}
% \tlabel[cc](0.5,4){\scriptsize $y$}
% \xmarks{-2,-1,1,2}
% \ymarks{1,2,3}
% \tlpointsep{4pt}
% \axislabels {x}{{\tiny $-2 \hspace{6pt}$} -2, {\tiny $-1 \hspace{6pt}$} -1, {\tiny $1$} 1, {\tiny $2$} 2}
% \axislabels {y}{{\tiny $1$} 1, {\tiny $2$} 2, {\tiny $3$} 3}
% \penwd{1.25pt}
% \arrow \reverse \arrow \function{-3,3,0.1}{2}
% \point[4pt]{(-2,2), (-1,2), (0,2), (1,2), (2,2)}
% \end{mfpic} 




\item   

$h(s) = s^3$ 

Domain: $(-\infty, \infty)$ 

Range:  $(-\infty, \infty)$


% \begin{mfpic}[10]{-3}{3}{-9}{9}

% \axes
% \tlabel[cc](3,-0.5){\scriptsize $s$}
% \tlabel[cc](0.5,9){\scriptsize $y$}
% \ymarks{-8,-7,-6,-5,-4,-3,-2,-1,1,2,3,4,5,6,7,8}
% \xmarks{-2,-1,1,2}
% \tlpointsep{4pt}
% \tiny 
% \axislabels {x}{{$-2 \hspace{6pt}$} -2, {$-1 \hspace{6pt}$} -1, {$1$} 1, {$2$} 2}
% \axislabels {y}{{$-8$} -8,{$-7$} -7,{$-6$} -6,{$-5$} -5,{$-4$} -4,{$-3$} -3,{$-2$} -2, {$-1$} -1, {$1$} 1, {$2$} 2, {$3$} 3, {$4$} 4, {$5$} 5, {$6$} 6, {$7$} 7, {$8$} 8}
% \normalsize
% \penwd{1.25pt}
% \arrow \reverse \arrow \parafcn{-2.1,2.1,0.1}{(t,t**3)}
% \point[4pt]{(0,0), (-1, -1), (1, 1), (-2, -8), (2, 8)}
% \end{mfpic}



\item  

$f(x) = x(x-1)(x+2)$

Domain: $(-\infty, \infty)$ 

Range: $(-\infty, \infty)$


 

 


% \begin{mfpic}[15]{-3}{3}{-1}{5}
% \axes
% \tlabel[cc](3,-0.5){\scriptsize $x$}
% \tlabel[cc](0.5,5){\scriptsize $y$}
% \xmarks{-2,-1,1,2}
% \ymarks{1,2,3,4}
% \tlpointsep{4pt}
% \axislabels {x}{{\tiny $-2 \hspace{6pt}$} -2, {\tiny $-1 \hspace{6pt}$} -1, {\tiny $1$} 1, {\tiny $2$} 2}
% \axislabels {y}{{\tiny $1$} 1, {\tiny $2$} 2, {\tiny $3$} 3, {\tiny $4$} 4}
% \penwd{1.25pt}
% \arrow \reverse \arrow \function{-2.25, 1.75, 0.1}{x*(x-1)*(x+2)}
% \point[4pt]{(-2,0), (-1,2), (0,0), (1,0)}
% \end{mfpic} 



\item

$g(t) = \sqrt{t-2}$

Domain: $[2, \infty)$ 

Domain: $[0, \infty)$ 


 

 


% \begin{mfpic}[15]{-1}{10}{-1}{4}
% \axes
% \tlabel[cc](10,-0.5){\scriptsize $t$}
% \tlabel[cc](0.5,3.75){\scriptsize $y$}
% \xmarks{1,2,3,4,5,6,7,8,9}
% \ymarks{1,2,3}
% \tlpointsep{4pt}
% \axislabels {x}{{\tiny $1$} 1, {\tiny $2$} 2, {\tiny $3$} 3, {\tiny $4$} 4, {\tiny $5$} 5, {\tiny $6$} 6, {\tiny $7$} 7, {\tiny $8$} 8, {\tiny $9$} 9}
% \axislabels {y}{{\tiny $1$} 1, {\tiny $2$} 2, {\tiny $3$} 3}
% \penwd{1.25pt}
% \arrow \function{2, 10, 0.1}{sqrt(x - 2)}
% \point[4pt]{(2,0), (3,1), (6,2)}
% \end{mfpic}



\item  

$h(s) = \sqrt{5 - s}$ 

Domain: $(-\infty, 5]$ 

Range:  $[0, \infty)$

% \begin{mfpic}[15]{-5}{6}{-1}{4}

% \axes
% \tlabel[cc](6,-0.5){\scriptsize $s$}
% \tlabel[cc](0.5,3.75){\scriptsize $y$}
% \xmarks{-4,-3,-2,-1,1,2,3,4,5}
% \ymarks{1,2,3}
% \tlpointsep{4pt}
% \tiny 
% \axislabels {x}{{$-4 \hspace{6pt}$} -4, {$-3 \hspace{6pt}$} -3, {$-2 \hspace{6pt}$} -2, {$-1 \hspace{6pt}$} -1, {$1$} 1, {$2$} 2, {$3$} 3, {$4$} 4, {$5$} 5}
% \axislabels {y}{{$1$} 1, {$2$} 2, {$3$} 3}
% \normalsize
% \penwd{1.25pt}
% \arrow \reverse \function{-4.5, 5, 0.1}{sqrt(5 - x)}
% \point[4pt]{(5,0), (0, 2.2360679), (1, 2), (-4, 3)}
% \end{mfpic}




\item   

$f(x) = 3-2\sqrt{x+2}$ 

Domain: $[-2,\infty)$ 

Range: $(-\infty, 3]$ 


 

 

% \begin{mfpic}[15]{-3}{3}{-1.5}{5}
% \axes
% \tlabel[cc](3,-0.5){\scriptsize $x$}
% \tlabel[cc](0.5,5){\scriptsize $y$}
% \xmarks{-2,-1,1,2}
% \ymarks{1,2,3,4}
% \tlpointsep{4pt}
% \axislabels {x}{{\tiny $-2 \hspace{6pt}$} -2, {\tiny $-1 \hspace{6pt}$} -1, {\tiny $1$} 1, {\tiny $2$} 2}
% \axislabels {y}{{\tiny $1$} 1, {\tiny $2$} 2, {\tiny $3$} 3, {\tiny $4$} 4}
% \penwd{1.25pt}
% \arrow \function{-2, 3, 0.1}{3-2*sqrt(x+2)}
% \point[4pt]{(-2,3), (-1,1), (2,-1)} 
% \end{mfpic} 




\item 

$g(t) = \sqrt[3]{t}$ 

Domain: $(-\infty, \infty)$ 

Range:  $(-\infty, \infty)$

 

% \begin{mfpic}[10]{-9}{9}{-3}{3}
% \point[4pt]{(0,0), (-1, -1), (1, 1), (-8, -2), (8, 2)}
% \axes
% \tlabel[cc](9,-0.5){\scriptsize $t$}
% \tlabel[cc](0.5,2.75){\scriptsize $y$}
% \xmarks{-8,-7,-6,-5,-4,-3,-2,-1,1,2,3,4,5,6,7,8}
% \ymarks{-2,-1,1,2}
% \tlpointsep{4pt}
% \tiny 
% \axislabels {x}{{$-8 \hspace{6pt}$} -8, {$-7 \hspace{6pt}$} -7, {$-6 \hspace{6pt}$} -6, {$-5 \hspace{6pt}$} -5, {$-4 \hspace{6pt}$} -4, {$-3 \hspace{6pt}$} -3, {$-2 \hspace{6pt}$} -2, {$-1 \hspace{6pt}$} -1, {$1$} 1, {$2$} 2, {$3$} 3, {$4$} 4, {$5$} 5, {$6$} 6, {$7$} 7, {$8$} 8}
% \axislabels {y}{{$-2$} -2, {$-1$} -1, {$1$} 1, {$2$} 2}
% \normalsize
% \penwd{1.25pt}
% \arrow \reverse \arrow \parafcn{-2.1,2.1,0.1}{(t**3,t)}
% \point[4pt]{(0,0), (-1, -1), (1, 1), (-8, -2), (8, 2)}
% \end{mfpic}



\item  

$h(s) = \dfrac{1}{s^{2} + 1}$ 

Domain: $(-\infty, \infty)$ 

Range:  $(0, 1]$

 

% \begin{mfpic}[23]{-3}{3}{-1}{2}

% \axes
% \tlabel[cc](3,-0.5){\scriptsize $s$}
% \tlabel[cc](0.5,1.75){\scriptsize $y$}
% \xmarks{-2,-1,1,2}
% \ymarks{1}
% \tlpointsep{4pt}
% \scriptsize
% \axislabels {x}{{$-2 \hspace{7pt}$} -2, {$-1 \hspace{7pt}$} -1, {$1$} 1, {$2$} 2}
% \axislabels {y}{{$1$} 1}
% \normalsize
% \penwd{1.25pt}
% \arrow \reverse \arrow \function{-2.5, 2.5, 0.1}{1/(x**2 + 1)}
% \point[4pt]{(0, 1), (1,0.5), (-1,0.5)}
% \end{mfpic}



%\setcounter{HW}{\value{enumi}}
\end{enumerate}


\begin{enumerate}
%\setcounter{enumi}{\value{HW}}

\item   

\begin{enumerate}

\item domain $ = \{ -1, 0, 1, 2 \}$, range $ = \{ -3, 0, 4\}$

\item  $f(0) = -3$,  $f(x) = 0$ for $x = -1, 1$.

\item  $f = \{ (-1,0), (0, -3), (1,0), (2,4) \}$

  

\item  $~$

\begin{image}
\begin{tikzpicture}[font=\LARGE]

\begin{axis}[
    xmin=-4, xmax=4,
    ymin=-5, ymax=5,
    width=5in,
    height=5in,
    xtick={-3,-2,-1,0,1,2,3},
    ytick={-4,-3,-2,-1,0,1,2,3,4},
    xticklabels={$-3$,$-2$,$-1$,$0$,$1$,$2$,$3$},
    yticklabels={$-4$,$-3$,$-2$,$-1$,$0$,$1$,$2$,$3$,$4$},
    tick style={line width=1.2pt},
    major tick length=6pt,
    axis x line=center,
    axis y line=center,
    x axis line style={-{Latex[length=10pt,width=8pt]}, line width=1.2pt},
    y axis line style={-{Latex[length=10pt,width=8pt]}, line width=1.2pt},
    xlabel={$x$},
    ylabel={$y$},
    every axis x label/.style={
        at={(current axis.right of origin)},
        anchor=west
    }
]

\addplot[
    only marks,
    mark=*,
    mark size=4pt
]
coordinates {
    (-1,0)
    (0,-3)
    (1,0)
    (2,4)
};

\end{axis}

\end{tikzpicture}
\end{image}

\end{enumerate}


\item    

\begin{enumerate}

\item  domain $= \{ -1, 0, 2, 3 \}$, range $=\{ 2, 3, 4 \}$

\item $~$

\begin{image}
\begin{tikzpicture}[font=\LARGE]

% Function name
\node at (0,6) {$g$};

% Domain values
\node at (-3,5) {$-1$};
\node at (-3,3) {$0$};
\node at (-3,1) {$2$};
\node at (-3,-1) {$3$};

% Codomain values
\node at (3.5,5) {$2$};
\node at (3.5,3) {$3$};
\node at (3.5,1) {$4$};

% Arrows
\draw[-{Latex[length=10pt,width=8pt]}, line width=1.5pt]
    (-2.5,5) -- (2.5,1);

\draw[-{Latex[length=10pt,width=8pt]}, line width=1.5pt]
    (-2.5,3) -- (2.5,5);

\draw[-{Latex[length=10pt,width=8pt]}, line width=1.5pt]
    (-2.5,1) -- (2.5,3);

\draw[-{Latex[length=10pt,width=8pt]}, line width=1.5pt]
    (-2.5,-1) -- (2.5,1);

\end{tikzpicture}
\end{image}
 

\item  Find $g(0) = 2$ and $g(x) = 0$ has no solutions.



\item $~$

\begin{image}
\begin{tikzpicture}[font=\LARGE]

\begin{axis}[
    xmin=-4, xmax=4,
    ymin=-1, ymax=5,
    width=5in,
    height=5in,
    xtick={-3,-2,-1,0,1,2,3},
    ytick={1,2,3,4},
    xticklabels={$-3$,$-2$,$-1$,$0$,$1$,$2$,$3$},
    yticklabels={$1$,$2$,$3$,$4$},
    tick style={line width=1.2pt},
    major tick length=6pt,
    axis x line=center,
    axis y line=center,
    x axis line style={-{Latex[length=10pt,width=8pt]}, line width=1.2pt},
    y axis line style={-{Latex[length=10pt,width=8pt]}, line width=1.2pt},
    xlabel={$x$},
    ylabel={$y$},
    every axis x label/.style={
        at={(current axis.right of origin)},
        anchor=west
    }
]

\addplot[
    only marks,
    mark=*,
    mark size=4pt
]
coordinates {
    (-1,4)
    (0,2)
    (2,3)
    (3,4)
};

\end{axis}

\end{tikzpicture}
\end{image}


\end{enumerate}


\item  $F(4) = 4^2 = 16$ (when $t = 4$), the solutions to $F(x) = 4$ are $x = \pm 2$ (when $t = \pm 2$). 

\item  $G(4) = 7$ (when $t = 2$), the solution to $G(t) = 4$ is $x = -2$ (when $t = -1$)


%\setcounter{HW}{\value{enumi}}

\end{enumerate}


\begin{enumerate}
%\setcounter{enumi}{\value{HW}}

\item  $A(3) = 9$, so the area enclosed by a square with a side of length $3$ inches is $9$ square inches.  The solutions to $A(\ell) = 36$ are $\ell = \pm 6$.  Since $\ell$ is restricted to  $\ell > 0$, we only keep $\ell  = 6$.  This means for the area enclosed by the square to be $36$ square inches, the length of the side needs to be $6$ inches.  Since $\ell$ represents a length, $\ell > 0$.



\item  $A(2) = 4\pi$, so the area enclosed by a circle with radius $2$ meters is $4\pi$ square meters.  The solutions to $A(r) = 16\pi$ are $r = \pm 4$.  Since $r$ is restricted to $r > 0$, we only keep $r = 4$.  This means for the area enclosed by the circle to be $16\pi$ square meters, the radius needs to be $4$ meters.  Since $r$ represents a radius (length), $r > 0$.

\item  $V(5) = 125$, so the volume enclosed by a cube with a side of length $5$ centimeters is $125$ cubic centimeters.  The solution to $V(s) = 27$ is $s = 3$.  This means for the volume enclosed by the cube to be $27$ cubic centimeters, the length of the side needs to $3$ centimeters.  Since $x$ represents a length, $x > 0$.

\item  $V(3) = 36\pi$, so the volume enclosed by a sphere with radius $3$ feet is $36\pi$ cubic feet.  The solution to $V(r) = \frac{32\pi}{3}$ is $r = 2$.  This means for the volume enclosed by the sphere to be $\frac{32\pi}{3}$ cubic feet, the radius needs to $2$ feet.  Since $r$ represents a radius (length), $r > 0$.


\item $h(0) = 64$, so at the moment the object is dropped off the building, the object is $64$ feet off of the ground.  The solutions to $h(t) = 0$ are $t = \pm 2$.  Since we restrict $0 \leq t \leq 2$, we only keep $t = 2$.  This means $2$ seconds after the object is dropped off the building, it is $0$ feet off the ground.  Said differently, the object hits the ground after $2$ seconds.  The restriction  $0 \leq t \leq 2$ restricts the time to be between the moment the object is released and the moment it hits the ground.


\item  $T(0) = 3$, so at 6 AM ($0$ hours after 6 AM), it is $3^{\circ}$ Fahrenheit.  $T(6) = 33$, so at noon ($6$ hours after 6 AM), the temperature is $33^{\circ}$ Fahrenheit.  $T(12) = 27$, so at 6 PM ($12$ hours after 6 AM), it is $27^{\circ}$ Fahrenheit.


\item $C(0) = 27$, so to make $0$ pens, it costs\footnote{This is called the `fixed' or `start-up' cost.  We'll revisit this concept in Example \ref{PortaBoyCost} in Section \ref{ConstantandLinearFunctions}.} $\$ 2700$.  $C(2) = 11$, so to make $2000$ pens, it costs $\$1100$.  $C(5) = 2$, so to make $5000$ pens, it costs $\$2000$.

\item $E(0) = 16.00$, so in 1980 ($0$ years after 1980), the average fuel economy of passenger cars in the US was $16.00$ miles per gallon.  $E(14) = 20.81$, so in 1994 ($14$ years after 1980), the average fuel economy of passenger cars in the US was $20.81$ miles per gallon.  $E(28) = 22.64$, so in 2008 ($28$ years after 1980), the average fuel economy of passenger cars in the US was $22.64$ miles per gallon.  


\item  $P(s) = 4s$, $s > 0$.

\item  $C(D) = \pi D$,  $D > 0$.

\item

\begin{enumerate}

\item The amount in the retirement account after 30 years if the monthly payment is $\$50$.

\item  The solution to $A(P) = 250000$ is what the monthly payment needs to be in order to have $\$250,\!000$ in the retirement account after 30 years.

\item  $A(P+50)$ is how much is in the retirement account in 30 years if $\$ 50$ is added to the monthly payment $P$.  $A(P)+50$ represents the amount of money in the retirement account after 30 years if $\$P$  is invested each month plus an additional $\$50$.  $A(P)+A(50)$ is the sum of money from two retirement accounts after 30 years: one with monthly payment $\$P$ and one with monthly payment $\$50$.

\end{enumerate}

\item  

\begin{enumerate}

\item  Since noon is $4$ hours after 8 AM, $P(4)$ gives the chance of precipitation at noon.

\item  We would need to solve $P(t) \geq 50 \%$ or $P(t) \geq 0.5$.

\end{enumerate}


\item The graph in question passes the horizontal line test meaning for each $w$ there is only one $v$.    The domain of $g$ is $[0, \infty)$ (which is the range of $f$) and the range of $g$ is $[2, \infty)$ which is the domain of $f$.  

\item  Answers vary.  

\end{enumerate}

\end{document}
